\documentclass[10pt,a4paper]{amsart}
\usepackage[margin=19mm]{geometry}
\usepackage{amsmath,amssymb,mathtools}
\usepackage[scr=rsfs,cal=euler]{mathalfa}
\usepackage{xfrac}
\usepackage[unicode,hidelinks,bookmarks=false]{hyperref}
\newcommand{\ie}{i.e.\ }
\newcommand{\cf}{cf.\ }
\newcommand{\resp}{respectively\ }
\newcommand{\Subsection}{\subsection}
\providecommand{\nobreakdash}{\nobreak\hbox{-}}
\setlength{\emergencystretch}{2em}
\multlinegap0pt
\usepackage{bm}
\usepackage{bbm}
\usepackage{dsfont}
\usepackage{xspace}
\usepackage{cancel}
\usepackage{mathtools}
\usepackage{amsthm}
\makeatletter
\@ifundefined{thm}{\newtheorem{thm}{Thm}}{}
\@ifundefined{cor}{\newtheorem{cor}[thm]{Corollary}}{}
\@ifundefined{defn}{\newtheorem{defn}[thm]{Definition}}{}
\@ifundefined{lem}{\newtheorem{lem}[thm]{Lemma}}{}
\@ifundefined{prop}{\newtheorem{prop}[thm]{Proposition}}{}
\makeatother
\title[On the pro-modularity in the residually reducible case for s]{On the pro-modularity in the residually reducible case for some totally real fields}
\author{Xinyao Zhang}
\date{Source: arXiv:2411.18661v4 (CC BY 4.0)}
\begin{document}
\maketitle

\noindent\emph{Source and licence.} On the pro-modularity in the residually reducible case for some totally real fields. Version arXiv:2411.18661v4 (CC BY 4.0), distributed under the Creative Commons Attribution 4.0 International licence, as recorded in the arXiv licence metadata for this exact version and shown on the abstract page https://arxiv.org/abs/2411.18661v4. Full rights notice: licenses/arxiv-2411.18661-CC-BY-4.0.txt.\par\smallskip

\noindent\textit{Editorial scope.} Counted unit: the Lem whose source label is \texttt{None} in the article (source0.tex lines 937--939, source label \texttt{None}), delivered together with the article's own complete original proof of it (source0.tex lines 941--989, one \texttt{proof} environment of 49 lines). No mathematics is written, repaired or re-derived: statement and proof are the article's own text line for line. Required context retained uncounted: pseudo (lines 212--236); u (lines 260--266); RX1 (lines 517--521); red-pseudo (lines 526--531); red-one (lines 539--545); nice (lines 682--701); equiv (lines 705--716); ordinary (lines 801--814); dense (lines 820--832). Every definition, display and label of those lines is retained unchanged. Omitted from this package and not redistributed: 1--211, 237--259, 267--516, 522--525, 532--538, 546--681, 702--704, 717--800, 815--819, 833--936, 940--940, 990--1329. Nothing inside a retained excerpt is omitted or altered: every retained body is delivered exactly as the article's lines are written, so no omission receipt of any kind accompanies this package. Citations: the retained excerpts carry no citation command at all, so no bibliography entry of the article is transcribed and no reference is redistributed. Reference adaptation: none was needed and none was made; every reference inside the delivered unit resolves inside it, so no unresolved reference or citation is produced. Printed slips kept literal: no printed word, symbol, hypothesis or number of the counted statement or of its proof was altered. Layout and class: the article's own class is \texttt{amsart} with packages including mathtools, titlesec, caption, subcaption, tikz-cd, amsmath, amsthm, amsfonts, amscd, amssymb, amsopn, enumerate, enumitem, mathrsfs; the wrapper is the lane's pinned \texttt{amsart} editorial wrapper at 10pt on A4, which restates the theorem-like environments the retained text needs and adds the article's own macro definitions that the retained text needs (none), with extra wrapper packages (bm, bbm, dsfont, xspace, cancel, mathtools). The delivered numbering is the unit's own. Assets: no retained span contains a figure environment, a graphics-inclusion command, a PDF-page inclusion, a table or a TikZ picture; no image file is copied into this package, the package declares no asset dependency and no image file is redistributed with it. Licence: version arXiv:2411.18661v4 of this article is distributed under the Creative Commons Attribution 4.0 International licence, as recorded in the arXiv licence metadata for this exact version and shown on the abstract page https://arxiv.org/abs/2411.18661v4. A full rights notice is delivered at licenses/arxiv-2411.18661-CC-BY-4.0.txt. Characters: every retained excerpt is pure ASCII, so the delivered engine, XeLaTeX with its default Latin Modern text font, reports no missing character.
\begin{defn} \label{pseudo}
	For a profinite group $G$ and a topological commutative ring $R$ in which 2 is invertible, a 2-dimensional \textit{pseudo-representation} is a continuous function $T: G \to R$ such that 
	
	1) $T(1)=2$,
	
	2) there exists an order 2 element $c \in G$ such that $T(c)=0$,
	
	3) $T(\sigma \tau)=T(\tau \sigma)$ for all $ \sigma, \tau \in G$,
	
	4) $T(\gamma \delta \eta)+T(\gamma \eta \delta)-T(\gamma \eta)T(\delta)-T(\eta \delta)T(\gamma)-T(\delta \gamma)T(\eta)+T(\gamma)T(\delta)T(\eta)=0,$
	for any $\delta, \gamma,\eta\in G$.
	
	The \textit{determinant} $\operatorname{det}(T)$ of a pseudo-representation $T$ is a character (using 4)) defined by $$ \operatorname{det}(T): G \to R^\times, ~
	\operatorname{det}(T)(\sigma)=\frac{1}{2}(T(\sigma)^2-T(\sigma^2)), \sigma \in G.$$
	
	
	
	For a pseudo-representation $T$ and $\sigma, \tau \in G$, define:
	
	1) $a(\sigma)=\frac{1}{2}(T(c\sigma)+T(\sigma))$,
	
	2) $d(\sigma)=\frac{1}{2}(-T(c\sigma)+T(\sigma))$,
	
	3) $y(\sigma, \tau)=a(\sigma \tau)-a(\sigma)a(\tau).$
\end{defn}

\begin{prop}\label{u}
	Let $R$ be an integral domain, and $T: G \to R$ be a pseudo-representation for a profinite group $G$. Let $\alpha, \beta $ be two elements in $ G$.
	
	1) If $y(\alpha, \beta) =0$ in $R$, then either $y(\alpha, \tau)=0$ or $y(\tau, \beta)=0$ for any $\tau \in G$.
	
	2) Assume that the set $J_{\beta} =\{\tau \in G : y(\beta, \tau)\ne 0\}$ (resp. $J'_{\beta} =\{\tau \in G : y(\tau, \beta) \ne 0\}$) is non-empty. For $\tau \in J_{\beta}$ (resp. $J'_{\beta}$), let $u(\alpha, \beta)(\tau)= \frac{y(\alpha, \tau)}{y(\beta, \tau)}$ (resp. $u'(\alpha, \beta)(\tau)=\frac{y(\tau, \alpha)}{y(\tau, \beta)}$) be an element in the fraction field of $R$. Then we have $u(\alpha, \beta)(\tau)=u(\alpha, \beta)(\tau')$ (resp. $u'(\alpha, \beta)(\tau)=u'(\alpha, \beta)(\tau')$)  for any $\tau, \tau' \in J_\beta$ (resp. $J'_{\beta}$). (Afterwards, we will use notations $u(\alpha, \beta)$ and $u'(\alpha, \beta)$ for simplicity.)
\end{prop}

\begin{prop}\label{RX1}
	1) $R_x^{\textnormal{aux}, \{\xi_v\}}/\pi=R_x^{\textnormal{aux}, \{\xi'_v\}}/\pi$.
	
	2) $R^{\textnormal{ps}, \{\xi_v\}}/\pi= R^{\textnormal{ps}, \{\xi'_v\}}/\pi$. 
\end{prop}

\begin{prop}\label{red-pseudo} Assume $F$ is abelian over $\mathbb{Q}$.
	
	1) We have $\operatorname{dim}(R^{\textnormal{ps,red}}) \le 2$.
	
	2) Assume $\bar{\chi}$ is quadratic and can be extended to $G_{\mathbb{Q}}$. Then the dihedral locus $C^{\operatorname{dih}}$ of $\operatorname{Spec} R_\textnormal{aux}^{\textnormal{ps}}$, i.e., for any point $\mathfrak{p}$ in $C^{\operatorname{dih}}$, $\rho(\mathfrak{p}) \cong \operatorname{Ind}_{G_L}^{G_F} \theta$ for some character $\theta$ and the splitting field $L=F(\bar{\chi})$, is of dimension at most $2+[F: \mathbb{Q}]$. 
\end{prop}

\begin{cor}\label{red-one}
	Keep assumptions as in Proposition \ref{red-pseudo}. Let $\mathfrak{s}$ be a prime of $R_\textnormal{aux}^\textnormal{ps} $ containing $p$ and write $R=R_\textnormal{aux}^\textnormal{ps}/\mathfrak{s}$.
	
	1) If $\dim R \ge 2$, then we can find a prime $\mathfrak{s}'$ of $R_\textnormal{aux}^\textnormal{ps} $ containing $\mathfrak{s}$ such that the corresponding representation is irreducible.
	
	2) If $\dim R \ge 2+[F: \mathbb{Q}]$, then we can find a prime $\mathfrak{s}'$ of $R_\textnormal{aux}^\textnormal{ps} $ containing $\mathfrak{s}$ such that the corresponding representation is irreducible and not dihedral.
\end{cor}

\begin{defn}\label{nice}
	(1) We say that a prime  of $ R^{\textnormal{ps}, \{\xi_v\}}$ is \textit{pro-modular} if it comes from a prime of $(\mathbb{T}_{\xi})_{\mathfrak{m}_{\xi}} $. We say that a prime of $ R_\textnormal{aux}^\textnormal{ps}$ is \textit{pro-modular} if it comes from a pro-modular prime of some $ R^{\textnormal{ps}, \{\xi_v\}}$.
	
	(2) Let $\mathfrak{q} $ be a prime of $ (\mathbb{T}_{\xi})_{\mathfrak{m}_{\xi}} $ and $A$ be the normal closure of $(\mathbb{T}_{\xi})_{\mathfrak{m}_{\xi}} / \mathfrak{q} $ in $k(\mathfrak{q})$. We say that $\mathfrak{q} $ is a \textit{nice} prime if $\mathfrak{q} $ contains $p$ and $\operatorname{dim} (\mathbb{T}_{\xi})_{\mathfrak{m}_{\xi}} / \mathfrak{q}=1$ and there exists a two-dimensional representation $\rho(\mathfrak{q})^o : \operatorname{Gal}(F_{\Sigma}/F) \to \operatorname{GL}_2(A)$ satisfying:
	
	\quad 1) $\rho(\mathfrak{q})^o \otimes k(\mathfrak{q}) \cong \rho(\mathfrak{q})$ is irreducible.
	
	\quad 2) The $\operatorname{mod}~\mathfrak{m}_A$ reduction $\bar{\rho}_b$ of $\rho(\mathfrak{q})^o$ is a non-split extension and has the form $ \bar{\rho}_b(g)=\begin{pmatrix}
		* & *\\
		0 & *
	\end{pmatrix}$, $g \in \operatorname{Gal}(F_{\Sigma}/F)$. Here $\mathfrak{m}_A $ is the maximal ideal of $A$.
	
	\quad 3) (dihedral condition) If $\rho(\mathfrak{q}) $ is dihedral, namely isomorphic to $\operatorname{Ind}_{G_L}^{G_F} \theta$ for some quadratic extension $L$ of $F$ and continuous character $\theta: G_L \to k(\mathfrak{q})^\times$, then $L \cap F(\zeta_p) = F$,
	where $\zeta_p $ is a primitive $p$-th root of unity.
	
	\quad 4) $\rho(\mathfrak{q})^o \mid_{G_{F_v}} = \bar{\rho}_b \mid_{G_{F_v}} $ for any $v \in \Sigma \setminus \Sigma_p$.
	
	(3) We say that a prime of $ R^{\textnormal{ps}, \{\xi_v\}}$ is \textit{nice} if it comes from a nice prime of $(\mathbb{T}_{\xi})_{\mathfrak{m}_{\xi}} $.
	
\end{defn}

\begin{prop}\label{equiv}
	Let $\mathfrak{q}$ be a prime ideal of $ R^{\textnormal{ps}, \{\xi_v\}}$ containing $p$ such that $ R^{\textnormal{ps}, \{\xi_v\}}/\mathfrak{q}$ is one-dimensional. Suppose $\rho(\mathfrak{q})$ is irreducible. Then the third condition of (2) in Definition \ref{nice} holds for $\mathfrak{q}$ if one of the following conditions holds:
	
	1)  $\bar{\chi}$ is not quadratic. 
	
	2) $\bar{\chi}\mid_{G_{F_v}}$ is trivial for some $v \mid p$.
	
	3) There exists a place $v \mid p$ such that $\bar{\chi}\mid_{G_{F_v}}$ is not trivial and $\rho(\mathfrak{q})^o\mid_{G_{F_v}} \cong \begin{pmatrix}
		\chi_{v,1} & *\\
		0 & \chi_{v,2}
	\end{pmatrix}$ is reducible. Moreover, $ \chi_{v,1}$ is of infinite order.
\end{prop}

\begin{prop}\label{ordinary}
	(1) $ R^{\textnormal{ps,ord}}$ is a finite $ \Lambda_F$-algebra. 
	
	(2) For any maximal ideal $\mathfrak{p}$ of $R^{\textnormal{ps,ord}}[\frac{1}{p}]$, we denote the associated semi-simple representation $\operatorname{Gal}(F_{\Sigma}/F) \to \operatorname{GL}_2(k(\mathfrak{p}))$ by $\rho(\mathfrak{p})$. Assume that
	
	\quad i) $ \rho(\mathfrak{p})$ is irreducible.
	
	\quad ii) For any $v \mid p$, $\rho(\mathfrak{p}) \mid_{G_{F_v}} \cong  \begin{pmatrix}
		\chi_{v,1} & *\\
		0 & \chi_{v,2}
	\end{pmatrix}$ such that $\chi_{v,1} $ is de Rham and has strictly less Hodge-Tate number than $\chi_{v,2} $ for any embedding $F_v \hookrightarrow \overline{\mathbb{Q}}_p$.
	
	Then $ \rho(\mathfrak{p})$ comes from a twist of Hilbert modular form.
\end{prop}

\begin{lem}\label{dense}
	Let $C^{\textnormal{ord}}$ be an irreducible component of $ R^{\textnormal{ps,ord}}$ of dimension at least $1+[F: \mathbb{Q}]$.
	 Let $ C^{\textnormal{ord,aut}}$ be the set of regular de Rham primes in $C^{\textnormal{ord}}$. Then $C^{\textnormal{ord,aut}} $ is dense in $C^{\textnormal{ord}}$. Here a prime $\mathfrak{q}$ is called regular de Rham prime if
	
	\quad i) $p \notin \mathfrak{q}$ and $R^{\textnormal{ps,ord}}/\mathfrak{q} $ is one-dimensional.
	
	\quad ii) The semi-simple representation $\rho(\mathfrak{q}): \operatorname{Gal}(F_{\Sigma}/F) \to \operatorname{GL}_2(k(\mathfrak{q}))$ (in the sense of Definition \ref{pseudo}) is irreducible.
	
	\quad iii) For any $v \mid p$, $\rho(\mathfrak{q}) \mid_{G_{F_v}} \cong  \begin{pmatrix}
		\chi_{v,1} & *\\
		0 & \chi_{v,2}
	\end{pmatrix}$ such that $\chi_{v,1} $ is de Rham and has strictly less Hodge-Tate number than $\chi_{v,2} $ for any embedding $F_v \hookrightarrow \overline{\mathbb{Q}}_p$.
\end{lem}

	\begin{lem}
		There exists a nice prime containing $\mathfrak{P}$.
	\end{lem}

	\begin{proof}
		By Proposition \ref{ordinary}, for any prime $\mathfrak{q}_1 \in  C_1^{\textnormal{ord,aut}}$, the subset of $C_1^{\textnormal{ord}} $ consisting of all regular de Rham primes as defined in Lemma \ref{dense}, we know that $\rho(\mathfrak{q}_1)$ comes from a twist of Hilbert modular form. By the density result in Lemma \ref{dense}, the image of any prime of $C_1^{\textnormal{ord}}$ in $\operatorname{Spec} R^{\textnormal{ps},\{\xi_v\}}$ is pro-modular.  
		
		By Corollary \ref{red-one}, we can choose an irreducible one-dimensional prime $\mathfrak{r}'$ of $R''$. We have shown that $\mathfrak{r}' $ is pro-modular. Now we claim that $\mathfrak{r}' $ is actually a nice prime in the sense of Definition \ref{nice}.
		
		Let $A$ be the normal closure of $R''/\mathfrak{r}'$. Then we obtain a continuous irreducible representation $\rho(\mathfrak{r}') : \operatorname{Gal}(F_{\Sigma}/F) \to \operatorname{GL}_2(A)$ associated to the pseudo-deformation in the following form $$
		\rho(\mathfrak{r}')(\sigma)=\begin{pmatrix}
			a(\sigma) & \frac{y(\sigma, \tau_{\mathfrak{r}'})}{y(\sigma_{\mathfrak{r}'}, \tau_{\mathfrak{r}'})}\\
			y(\sigma_{\mathfrak{r}'}, \sigma)& d(\sigma)
	\end{pmatrix}, ~y(\sigma_{\mathfrak{r}'}, \tau_{\mathfrak{r}'}) \ne 0.$$
		
		Now we verify the conditions in Definition \ref{nice} for all the three cases. 
		
		In case a), we can take $\rho(\mathfrak{r}')^o=\rho(\mathfrak{r}')$. From our construction, the first three conditions are clear (using 3) of Proposition \ref{equiv}). By Proposition \ref{RX1}, we know that if $a(\sigma_{v})=1$, then we have $d(\sigma_{v})=1$ for any $v \in \Sigma \setminus \Sigma_p$ as $T(\sigma_{v})=\xi_v+\xi_v^{-1} \equiv 2~ \operatorname{mod}~\pi$. Thus, we can find that for all $v \in \Sigma \setminus \Sigma_p$, $\rho(\mathfrak{r}')^o \mid_{G_{F_v}}= \begin{pmatrix}
			1 & 0\\
			0 & 1
		\end{pmatrix}$. Then $\mathfrak{r}' $ is a nice prime.
		
		In case b), if $y(\theta_{i}, \alpha)=0$ for all $\alpha \in \operatorname{Gal}(F_\Sigma/F)$, then it is the same as case a). If not, by Proposition \ref{u}, we know that $y(\theta_{i}, \tau_{\mathfrak{r}'}) \ne  0$, otherwise $y(\sigma_{\mathfrak{r}'}, \tau_{\mathfrak{r}'}) = 0$. By 3) in \textit{Step} I, $u(\sigma_{\mathfrak{r}'}, \theta_{i})$ is well-defined and integral over $R''/\mathfrak{r}'$, and hence a unit in $A$. Then we can take $\rho(\mathfrak{r}')^o=P\rho(\mathfrak{r}')P^{-1}$, where $$P= \begin{pmatrix}
			u(\sigma_{\mathfrak{r}'}, \theta_{i}) & 0\\
			0 & 1
		\end{pmatrix} \in \operatorname{GL}_2(A).$$
		
		Similar to the case a), we only need to check the fourth condition. From our construction, we know that either $u(\theta, \sigma_{\mathfrak{r}'})=0$ or $u(\theta, \theta_i)^{n}=1$ for any $\theta \in S$. As $(n,p)=1$, by Hensel's lemma, $ u(\theta, \theta_{i})$ is an element of $\mathbb{F}' \hookrightarrow \mathbb{F}'[[T]]=A$, where $\mathbb{F}' $ is the residue field of $A$. Thus, the fourth condition holds.
		
		In case c), we consider the continuous irreducible Galois representation $\rho(\mathfrak{r}')' : \operatorname{Gal}(F_{\Sigma}/F) \to \operatorname{GL}_2(A)$ in the following form $$
		\rho(\mathfrak{r}')'(\sigma)=\begin{pmatrix}
			a(\sigma) & y(\sigma, \tau_{\mathfrak{r}'})\\
			\frac{y(\sigma_{\mathfrak{r}'}, \sigma)}{y(\sigma_{\mathfrak{r}'}, \tau_{\mathfrak{r}'})}& d(\sigma)
		\end{pmatrix}, ~y(\sigma_{\mathfrak{r}'}, \tau_{\mathfrak{r}'}) \ne 0.$$
		
		If $y(\alpha, \theta_j)=0$ for all $\alpha \in \operatorname{Gal}(F_\Sigma/F)$, then we can take $$
		\rho(\mathfrak{r}')^o=\begin{pmatrix}
			0 & 1\\
			1 & 0
		\end{pmatrix} \rho(\mathfrak{r}')'\begin{pmatrix}
		0 & 1\\
		1 & 0
		\end{pmatrix} \subset \operatorname{GL}_2(A). $$
		Similar to case a), we can vefiry that $\mathfrak{r}'$ is a nice prime.
		If not, we can take $$ \rho(\mathfrak{r}')^o=\begin{pmatrix}
			0 & u'(\tau_{\mathfrak{r}'}, \theta_{j})\\
			1 & 0
		\end{pmatrix} \rho(\mathfrak{r}')'\begin{pmatrix}
			0 & 1\\
			u'(\theta_{j}, \tau_{\mathfrak{r}'}) & 0
		\end{pmatrix} \subset \operatorname{GL}_2(A). $$
	Similar to case b), we can also vefiry that $\mathfrak{r}'$ is a nice prime.
	\end{proof}

\end{document}
